% GENERATED FILE. Edit canonical lessons, then run npm run generate:books.
% canonical-source-hash: fnv1a64:a26fc520c82bbe47
% canonical-lessons: HI-C553-siva, HI-C553-silsila, HI-C553-kram, HI-C553-caran, HI-C553-tithi

\chapter{In the Right Order}
\label{ch:hi-in-the-right-order}
\hlchaptermodality{553}

\begin{chapteropening}
\textbf{By the end of this chapter:} \emph{I can talk about order, stages, dates and except.}

Complete the last lesson of chapter 553: I can talk about order, stages, dates and except.
\end{chapteropening}

\section[\texorpdfstring{\hi{सिवा} (sivā)}{सिवा (sivā)}]{\hi{सिवा} (sivā) — except (\hi{के} \hi{सिवा})}

\label{lesson:HI-C553-siva}



\begin{quote}
Before the new one: say the Hindi for well, then the Hindi for but (in a lighter register).
\end{quote}

\subsection*{You'll want to know: \hi{सिवा}}
\textbf{\hi{सिवा}} — \emph{sivā} — \textquotedblleft{}except (\hi{के} \hi{सिवा})\textquotedblright{}.

\begin{quote}
\hi{रविवार} \hi{के} \hi{सिवा} \hi{दुकान} \hi{हर} \hi{दिन} \hi{खुली} \hi{रहती} \hi{है।}
\end{quote}

\begin{quote}
\emph{Ravivār ke sivā dukān har din khulī rahtī hai.}
\end{quote}

\begin{quote}
The shop is open every day except Sunday.
\end{quote}

\begin{cousinweb}
\textbf{\hi{सिवा}} (\emph{sivā}) is usually traced to Arabic \emph{siwā}, \textquotedblleft{}except\textquotedblright{}, through Persian.
\end{cousinweb}

\begin{grammarlens}[title={What you've built}]
One more word for saying what happens if.
\end{grammarlens}

\subsection*{Your turn}
\begin{itemize}
\raggedright
  \item \emph{Say it:} \emph{sivā}
  \item \emph{Say it:} \emph{sivā}, once more
  \item \emph{Say it:} \emph{Ravivār ke sivā dukān har din khulī rahtī hai.} — The shop is open every day except Sunday.
  \item \emph{Recall:} say the Hindi for well, then the Hindi for but (in a lighter register), then say \emph{sivā} again
\end{itemize}

\begin{tcolorbox}[breakable,colback=teal!4,colframe=teal!35!black,title={Before you move on}]
What does \hi{सिवा} mean? (\textquotedblleft{}Except (\hi{के} \hi{सिवा})\textquotedblright{}.) Say it once more.
\end{tcolorbox}
\section[\texorpdfstring{\hi{सिलसिला} (silsilā)}{सिलसिला (silsilā)}]{\hi{सिलसिला} (silsilā) — a series; a run (of events)}

\label{lesson:HI-C553-silsila}



\begin{quote}
Before the new one: say the Hindi for but (in a lighter register), then the Hindi for except (\hi{के} \hi{सिवा}).
\end{quote}

\subsection*{You'll want to know: \hi{सिलसिला}}
\textbf{\hi{सिलसिला}} — \emph{silsilā} — \textquotedblleft{}a series; a run (of events)\textquotedblright{}.

\begin{quote}
\hi{बारिश} \hi{का} \hi{सिलसिला} \hi{तीन} \hi{दिन} \hi{से} \hi{चल} \hi{रहा} \hi{है।}
\end{quote}

\begin{quote}
\emph{Bāriś kā silsilā tīn din se cal rahā hai.}
\end{quote}

\begin{quote}
The rain has gone on for three days in a row.
\end{quote}

Masculine.

\begin{cousinweb}
\textbf{\hi{सिलसिला}} (\emph{silsilā}) is usually traced to Arabic \emph{silsila}, \textquotedblleft{}a chain\textquotedblright{}, through Persian.
\end{cousinweb}

\begin{grammarlens}[title={What you've built}]
One more word for saying what happens if.
\end{grammarlens}

\subsection*{Your turn}
\begin{itemize}
\raggedright
  \item \emph{Say it:} \emph{silsilā}
  \item \emph{Say it:} \emph{silsilā}, once more
  \item \emph{Say it:} \emph{Bāriś kā silsilā tīn din se cal rahā hai.} — The rain has gone on for three days in a row.
  \item \emph{Recall:} say the Hindi for but (in a lighter register), then the Hindi for except (\hi{के} \hi{सिवा}), then say \emph{silsilā} again
\end{itemize}

\begin{tcolorbox}[breakable,colback=teal!4,colframe=teal!35!black,title={Before you move on}]
What does \hi{सिलसिला} mean? (\textquotedblleft{}A series; a run (of events)\textquotedblright{}.) Say it once more.
\end{tcolorbox}
\section[\texorpdfstring{\hi{क्रम} (kram)}{क्रम (kram)}]{\hi{क्रम} (kram) — order; sequence}

\label{lesson:HI-C553-kram}



\begin{quote}
Before the new one: say the Hindi for except (\hi{के} \hi{सिवा}), then the Hindi for a series.
\end{quote}

\subsection*{You'll want to know: \hi{क्रम}}
\textbf{\hi{क्रम}} — \emph{kram} — \textquotedblleft{}order; sequence\textquotedblright{}.

\begin{quote}
\hi{शब्दों} \hi{का} \hi{क्रम} \hi{बदलिए।}
\end{quote}

\begin{quote}
\emph{Śabdoṁ kā kram badliye.}
\end{quote}

\begin{quote}
Change the order of the words.
\end{quote}

Masculine.

\begin{cousinweb}
\textbf{\hi{क्रम}} (\emph{kram}) is a tatsama, taken whole from Sanskrit \emph{krama}, \textquotedblleft{}a step; order, succession\textquotedblright{} (Monier-Williams).
\end{cousinweb}

\begin{grammarlens}[title={What you've built}]
One more word for saying what happens if.
\end{grammarlens}

\subsection*{Your turn}
\begin{itemize}
\raggedright
  \item \emph{Say it:} \emph{kram}
  \item \emph{Say it:} \emph{kram}, once more
  \item \emph{Say it:} \emph{Śabdoṁ kā kram badliye.} — Change the order of the words.
  \item \emph{Recall:} say the Hindi for except (\hi{के} \hi{सिवा}), then the Hindi for a series, then say \emph{kram} again
\end{itemize}

\begin{tcolorbox}[breakable,colback=teal!4,colframe=teal!35!black,title={Before you move on}]
What does \hi{क्रम} mean? (\textquotedblleft{}Order; sequence\textquotedblright{}.) Say it once more.
\end{tcolorbox}
\section[\texorpdfstring{\hi{चरण} (caraṇ)}{चरण (caraṇ)}]{\hi{चरण} (caraṇ) — a stage, a phase}

\label{lesson:HI-C553-caran}



\begin{quote}
Before the new one: say the Hindi for a series, then the Hindi for order.
\end{quote}

\subsection*{You'll want to know: \hi{चरण}}
\textbf{\hi{चरण}} — \emph{caraṇ} — \textquotedblleft{}a stage, a phase\textquotedblright{}.

\begin{quote}
\hi{योजना} \hi{का} \hi{पहला} \hi{चरण} \hi{पूरा} \hi{हो} \hi{गया।}
\end{quote}

\begin{quote}
\emph{Yojnā kā pahlā caraṇ pūrā ho gayā.}
\end{quote}

\begin{quote}
The first stage of the plan is complete.
\end{quote}

Masculine.

\begin{cousinweb}
\textbf{\hi{चरण}} (\emph{caraṇ}) is a tatsama, taken whole from Sanskrit \emph{caraṇa}, \textquotedblleft{}a foot; a step\textquotedblright{} (Monier-Williams). A step of a plan is a stage of it.
\end{cousinweb}

\begin{grammarlens}[title={What you've built}]
One more word for saying what happens if.
\end{grammarlens}

\subsection*{Your turn}
\begin{itemize}
\raggedright
  \item \emph{Say it:} \emph{caraṇ}
  \item \emph{Say it:} \emph{caraṇ}, once more
  \item \emph{Say it:} \emph{Yojnā kā pahlā caraṇ pūrā ho gayā.} — The first stage of the plan is complete.
  \item \emph{Recall:} say the Hindi for a series, then the Hindi for order, then say \emph{caraṇ} again
\end{itemize}

\begin{tcolorbox}[breakable,colback=teal!4,colframe=teal!35!black,title={Before you move on}]
What does \hi{चरण} mean? (\textquotedblleft{}A stage, a phase\textquotedblright{}.) Say it once more.
\end{tcolorbox}
\section[\texorpdfstring{\hi{तिथि} (tithi)}{तिथि (tithi)}]{\hi{तिथि} (tithi) — a date (formal); a day of the lunar month}

\label{lesson:HI-C553-tithi}



\begin{quote}
Before the new one: say the Hindi for order, then the Hindi for a stage, a phase.
\end{quote}

\subsection*{You'll want to know: \hi{तिथि}}
\textbf{\hi{तिथि}} — \emph{tithi} — \textquotedblleft{}a date (formal); a day of the lunar month\textquotedblright{}.

\begin{quote}
\hi{परीक्षा} \hi{की} \hi{तिथि} \hi{बदल} \hi{गई} \hi{है।}
\end{quote}

\begin{quote}
\emph{Parīkṣā kī tithi badal gaī hai.}
\end{quote}

\begin{quote}
The date of the exam has changed.
\end{quote}

Feminine.

\begin{cousinweb}
\textbf{\hi{तिथि}} (\emph{tithi}) is a tatsama, taken whole from Sanskrit \emph{tithi}, \textquotedblleft{}a lunar day\textquotedblright{} (Monier-Williams).
\end{cousinweb}

\begin{grammarlens}[title={What you've built}]
One more word for saying what happens if.
\end{grammarlens}

\subsection*{Your turn}
\begin{itemize}
\raggedright
  \item \emph{Say it:} \emph{tithi}
  \item \emph{Say it:} \emph{tithi}, once more
  \item \emph{Say it:} \emph{Parīkṣā kī tithi badal gaī hai.} — The date of the exam has changed.
  \item \emph{Recall:} say the Hindi for order, then the Hindi for a stage, a phase, then say \emph{tithi} again
\end{itemize}

\begin{tcolorbox}[breakable,colback=teal!4,colframe=teal!35!black,title={Before you move on}]
What does \hi{तिथि} mean? (\textquotedblleft{}A date (formal); a day of the lunar month\textquotedblright{}.) Say it once more.
\end{tcolorbox}
